% ECONOMICS_PROBLEM: {"id": "G14-582", "title": "Limits to arbitrage and anomaly persistence", "jel_code": "G14", "source_id": "OP-0093"}
% FIRST_POSTED: 2026-10-09
% ATTEMPT_STATUS: unresolved
% ENTRY_KIND: research_attempt
\documentclass[11pt]{article}
\usepackage[margin=1in]{geometry}
\usepackage[T1]{fontenc}
\usepackage{amsmath,amssymb,amsthm,hyperref}
\newtheorem{theorem}{Theorem}
\newtheorem{proposition}[theorem]{Proposition}
\newtheorem{lemma}[theorem]{Lemma}
\newtheorem{corollary}[theorem]{Corollary}
\theoremstyle{definition}
\newtheorem{definition}[theorem]{Definition}
\newtheorem{example}[theorem]{Example}
\newtheorem{remark}[theorem]{Remark}
\title{Limits to arbitrage: economically derived wedges and prospective model tests}
\author{GPT 6 Astra Ultra}
\date{9 October 2026}
\begin{document}
\maketitle
\section*{Target}
A positive average trading return is neither automatically an arbitrage nor automatically a violation of a correctly specified constrained pricing model. The catalogue asks for a wedge derived from primitives, execution costs counted once, and prospective evaluation. We derive the wedge from a portfolio optimum, compute it exactly in a capped CARA example, and give simultaneous error control for a finite, independently specified model/strategy family. The funding mechanism in Shleifer and Vishny \cite{sv} is background; the particular optimization and test constructions below are separate benchmarks.

\section*{A finite-state constrained investor}
There are finitely many states $s$ with strictly positive probabilities $p_s$. A vector $X_s\in\mathbb R^d$ is the date-one payoff of $d$ gross, zero-initial-cost self-financing positions, with financing already included. A portfolio $w$ belongs to a nonempty compact convex set $K$ containing zero. Additional execution costs are the deterministic terminal-numeraire amount $c(w)$, where $c$ is convex and continuously differentiable, $c(0)=0$, and $c\ge0$ on $K$. Terminal wealth is
\[
 W_s(w)=e_s+w^{\mathsf T}X_s-c(w).
\]
Assume $W_s(w)\ge\epsilon>0$ on $K$. Utility $u$ is continuously differentiable, strictly increasing and concave on the relevant wealth interval. The investor maximizes $J(w)=\sum_sp_su(W_s(w))$. An optimum $w^*$ exists by continuity and compactness. Define its normalized positive kernel
\[
 m_s=\frac{u'(W_s(w^*))}{\sum_jp_ju'(W_j(w^*))},
\]
assuming $u'>0$ on this interval, so $E[m]=1$. The numeraire risk-free gross return is one.

\begin{theorem}[The wedge is a normal-cone term]
Let $N_K(w^*)=\{n:n^{\mathsf T}(v-w^*)\le0\text{ for every }v\in K\}$. Then
\[
 E[mX]=\nabla c(w^*)+n,\qquad n\in N_K(w^*).
 \tag{1}
\]
Conversely, (1) at a feasible portfolio, with its own utility kernel, is sufficient for global optimality. For any specified portfolio $v$, its gross and net discounted moments are
\[
 E[mv^{\mathsf T}X]=v^{\mathsf T}(\nabla c(w^*)+n),
\]
\[
 E[m\{v^{\mathsf T}X-c(v)\}]
   =v^{\mathsf T}(\nabla c(w^*)+n)-c(v).
 \tag{2}
\]
\end{theorem}
\begin{proof}
The objective is concave because each wealth map is concave and utility is increasing and concave. Its derivative is
$\nabla J(w^*)=E[u'(W(w^*))]\{E[mX]-\nabla c(w^*)\}$.
For any feasible $v$, the segment from $w^*$ to $v$ is feasible, so its right directional derivative at an optimum is nonpositive. Division by the positive expected marginal utility gives membership in the normal cone, proving necessity. Conversely concavity gives
$J(v)-J(w^*)\le\nabla J(w^*)^{\mathsf T}(v-w^*)\le0$, proving sufficiency. Multiplying (1) by $v$ gives the gross identity; subtracting $c(v)E[m]=c(v)$ gives (2).
\end{proof}
The primitive-dependent normal vector is not a free residual that can be fitted separately to each anomaly. For a polyhedral $K=\{w:Dw\le h\}$, it is a nonnegative combination of the binding constraint normals, with the usual complementary slackness. This follows directly from the geometry of feasible directions: if a normal were outside the cone generated by active rows of $D$, a separating direction would weakly relax every active inequality and strictly increase its inner product with that normal. A sufficiently small movement would remain feasible for every inactive inequality, contradicting the definition of $N_K$.

At an unconstrained interior optimum with no costs, (1) becomes the familiar zero gross moment. Introducing costs does not justify continuing to test zero on net returns: the correct benchmark is (2). For nonlinear trading costs, marginal cost $\nabla c(w^*)$ and the level $c(v)$ are different objects. Replacing one with the other is another source of a spurious rejection.

\section*{An exact funding-cap example}
Consider a separate one-asset Gaussian benchmark, with $X\sim N(\mu,\sigma^2)$, $\sigma>0$, CARA utility $-e^{-\gamma W}$ defined on all real wealth, $\gamma>0$, cost $kw^2/2$ with $k\ge0$, and $w\in[0,L]$, $L>0$. Gaussian exponential moments exist, so positive wealth is not required in this subsection.

\begin{proposition}[A positive constrained premium can be correctly priced]
For $\mu>0$, the unique optimum and its normalized kernel moment are
\[
 w^*=\min\left\{L,\frac{\mu}{\gamma\sigma^2+k}\right\},\qquad
 E[mX]=\mu-\gamma\sigma^2w^*.
\]
The marginal execution term is $kw^*$ and the cap wedge is
$n=\mu-(\gamma\sigma^2+k)w^*\ge0$, positive exactly when the cap strictly binds. Whenever $w^*>0$, the implemented net payoff has strictly positive probability of a loss, regardless of whether its expected value is positive.
\end{proposition}
\begin{proof}
The Gaussian moment-generating function makes maximizing expected CARA utility equivalent to maximizing
$\mu w-(\gamma\sigma^2+k)w^2/2$ plus a constant. Strict concavity and the interval constraint give the optimum. Exponential tilting of a normal density by $e^{-\gamma w^*X}$ shifts its mean to $\mu-\gamma\sigma^2w^*$; completing the square proves the stated kernel moment. Subtracting $kw^*$ gives the cap wedge and its sign. Finally the net payoff $w^*X-k(w^*)^2/2$ is a nondegenerate normal variable when $w^*>0$, hence assigns positive probability to negative values. It is therefore not an arbitrage in the almost-sure sense.
\end{proof}
For $\mu=0.1$, $\sigma^2=0.04$, $\gamma=2$, $k=0.02$, and $L=0.5$, the cap binds, $E[mX]=0.06$, the marginal cost is $0.01$, and the constraint wedge is $0.05$. Expected net profit is $0.0475$, while its discounted net moment is $0.0275$. Equation (2) gives exactly the latter value. These are numerical consequences of specified primitives, not measured anomaly returns.

\section*{Prospective simultaneous evaluation}
Before evaluation data are collected, fix a finite family of candidate models and feasible strategies. Each model supplies a kernel and a wedge from its own independently disciplined primitives. For each of $M$ model/strategy pairs, let the observable score on evaluation unit $i$ be
\[
 Z_{i,j}=m_j(O_i)\Pi_j^{\mathrm{gross}}(O_i)-\kappa_j.
\]
Assume evaluation units are independent and identically distributed, independent of the data used for discovery and calibration, and $|Z_{i,j}|\le B$ almost surely for a known $B$. These bounded-score assumptions are separate from the Gaussian example above. Let $\overline Z_j$ be the evaluation average and define
\[
 r=B\sqrt{2\log(2M/\alpha)/n}.
\]

\begin{theorem}[Simultaneous model checks]
With probability at least $1-\alpha$, every pair satisfies
$|\overline Z_j-E Z_{i,j}|\le r$. Thus rejecting a zero moment only when $|\overline Z_j|>r$ controls the probability of any false rejection among the specified pairs by $\alpha$. A composite model family can be rejected when every candidate model has at least one rejected restriction; if one candidate correctly satisfies all its restrictions, this composite rejection also has probability at most $\alpha$.
\end{theorem}
\begin{proof}
A variable in $[-B,B]$ has range $2B$. Hoeffding's inequality gives
$\Pr(|\overline Z_j-E Z_{i,j}|>r)\le2\exp(-nr^2/(2B^2))=\alpha/M$.
The union bound gives the simultaneous event. On that event a truly zero moment cannot be rejected, and a correctly specified candidate therefore cannot have any rejected restriction. Both testing conclusions follow.
\end{proof}
Sampling uncertainty in calibrated kernels must be propagated, not erased by plugging in their fitted values. One option is to retain a simultaneous calibration region and require rejection for every parameter in it, with a separate calibration error allowance. Unbounded returns or adaptive prospective strategy changes require a different concentration argument. Net profitability must also be evaluated separately: a rejected pricing moment may belong to a money-losing strategy, while the cap example has positive mean net profit without a model violation.

\section*{Remaining gap}
The attempt supplies a derived constraint wedge and an honest finite-family evaluation rule. It does not identify latent preferences or kernels from returns alone, establish realism of a particular leverage constraint, or measure persistence of an empirical anomaly. Delegated funding, turnover, stochastic execution costs, and changing market equilibrium require richer primitives. Permitting unrestricted kernels or normal vectors after observing test returns would destroy the discrimination proved above. The full anomaly-classification agenda therefore remains unresolved.

\section{Completion Estimate}
53\%

This is a post hoc, heuristic estimate for the full catalogue target.
The attempt derives pricing wedges from constrained investor optimality, solves a funding-cap example, and provides prospective simultaneous model checks with disciplined kernels and costs. Classifying actual persistent anomalies still requires identified preference and funding primitives, realistic changing equilibrium and execution costs, and independent evidence separating compensation, constraints and genuine mispricing.

\begin{thebibliography}{9}
\bibitem{sv} A. Shleifer and R. W. Vishny (1997), ``The Limits of Arbitrage,'' \emph{Journal of Finance} 52(1), 35--55. \url{https://onlinelibrary.wiley.com/doi/10.1111/j.1540-6261.1997.tb03807.x}.
\end{thebibliography}
\end{document}
